\section{Overview}
The design of VeriGate draws on several distinct research threads: the transformer models that make semantic matching possible, the semantic-caching systems that VeriGate improves upon, the sentence-embedding and vector-search techniques that make a cache lookup fast, the Natural Language Inference (NLI) literature that motivates the verifier, and the emerging body of work on cost-aware LLM routing that underpins the second contribution. This chapter reviews each in turn (Sections 2.2--2.6) and closes with a gap analysis that situates VeriGate against the closest existing systems (Section~2.7).

\section{Large Language Models and Efficient Inference}
Modern LLMs are built on the transformer architecture \cite{vaswani2017attention}, whose self-attention mechanism enabled the pre-training of deep bidirectional \cite{devlin2019bert} and autoregressive language models at scale. Their capability, however, comes with substantial per-request cost and latency, and a growing survey literature catalogues techniques for making inference cheaper: quantisation, distillation, batching, KV-cache reuse, and request-level caching \cite{zhu2023survey}. VeriGate operates at the outermost of these layers --- the request level --- where savings are model-agnostic and require no access to model weights. This makes the approach portable across providers, which is essential for a gateway that must front closed, hosted APIs such as OpenAI and Gemini.

\section{Semantic Caching for LLM Applications}
The idea of caching LLM responses by \emph{meaning} rather than by exact string was crystallised by GPTCache \cite{bang2023gptcache}, which embeds each query, stores the embedding alongside the answer, and serves a cached answer when a new query's nearest neighbour exceeds a configured similarity threshold. GPTCache demonstrated meaningful cost and latency reductions and has become the reference open-source implementation of the pattern.

The critical limitation, acknowledged in practice, is the reliance on a \textbf{single fixed similarity threshold} together with the implicit assumption that \emph{high similarity implies semantic equivalence}. This assumption fails precisely in the cases that matter: two queries that differ by a single number, a negation, or a named entity are lexically and geometrically close yet demand different answers. A fixed threshold cannot resolve this, because lowering it to admit legitimate paraphrases simultaneously admits these look-alikes, while raising it to exclude look-alikes also rejects legitimate paraphrases. VeriGate's thesis is that the threshold decision and the equivalence decision must be \emph{separated}: the threshold proposes a candidate, and an independent verifier disposes of it. This separation is the primary conceptual departure from GPTCache-style systems.

\section{Text Embeddings and Semantic Similarity}
The quality of a semantic cache is bounded by the quality of its embeddings. Early text representations used sparse term-weighting such as TF-IDF \cite{salton1988term}, which captures lexical overlap but not meaning; two paraphrases with no shared content words are orthogonal under TF-IDF. Dense sentence embeddings solve this. Sentence-BERT \cite{reimers2019sentencebert} fine-tunes transformer encoders in a siamese configuration so that cosine similarity in the embedding space reflects semantic similarity, enabling efficient nearest-neighbour matching over sentences. VeriGate uses \texttt{all-MiniLM-L6-v2} \cite{sentencetransformers}, a distilled six-layer encoder \cite{wang2020minilm} producing 384-dimensional vectors that balances quality against the low inference latency a cache demands. A practical finding relevant to this literature, reported in Chapter~6, is that the similarity distribution --- and therefore the correct base threshold --- is a property of the specific embedding model, which is itself an argument against any universally fixed threshold.

\section{Approximate Nearest Neighbour Search and Vector Indexing}
For the cache lookup to be worthwhile it must be far faster than the LLM call it replaces, and it must remain fast as the cache grows. A brute-force scan over stored vectors is $O(n)$ per query. The information-retrieval community has produced sub-linear alternatives: product quantisation \cite{jegou2011pq} and GPU-accelerated exhaustive search \cite{johnson2019faiss} for compression and throughput, and graph-based indexes such as Hierarchical Navigable Small World (HNSW) \cite{malkov2018hnsw} for logarithmic-time approximate search. Production data stores now expose these directly; Redis Stack's vector-search capability builds an HNSW index over stored vectors with metadata filtering \cite{redis_stack}. VeriGate adopts a two-backend design informed by this literature: a vectorised in-process index (a single matrix product) for single-node deployments, and an optional RediSearch/HNSW backend for horizontally-scaled deployments that share one cache across replicas.

\section{Natural Language Inference and Answer Verification}
The verifier that distinguishes VeriGate from prior caches is grounded in Natural Language Inference (NLI), the task of deciding whether one sentence entails, contradicts, or is neutral toward another. The task was scaled by large annotated corpora such as SNLI \cite{bowman2015snli}, and modern cross-encoders based on DeBERTa \cite{he2021deberta} achieve strong entailment accuracy. VeriGate's Tier-2 verifier uses a compact NLI cross-encoder to test \emph{bidirectional} entailment: a cached answer is served only if each query entails the other. Because NLI cross-encoders are comparatively expensive, VeriGate precedes them with a Tier-1 structural check --- inspired by the observation that most dangerous false hits differ in a small set of high-signal features (numbers, negations, named entities) that can be compared almost for free. The two-tier design, and the empirical finding that Tier-2 is largely redundant once the adaptive threshold is active, are contributions of this work reported in Chapter~6.

\section{Cost-Aware Routing and Model Cascades}
The second contribution builds on recent work showing that not every query needs the most expensive model. FrugalGPT \cite{chen2023frugalgpt} demonstrated that cascading from cheap to expensive models, with an early-exit decision, can match the quality of a single large model at a fraction of the cost. RouteLLM \cite{ong2024routellm} learns a router from human-preference data to direct queries between a weak and a strong model. VeriGate's cost-aware cascade is a lightweight, dependency-free realisation of this idea: a transparent complexity score routes easy queries to a cheap model and escalates only hard ones, with the unused tier retained as a failover. It is deliberately interpretable rather than learned, so its routing decisions are auditable in a viva setting; learned routing in the style of RouteLLM is noted as future work.

\section{Gap Analysis and VeriGate's Contribution}
The reviewed literature reveals a clear gap. Semantic caching \cite{bang2023gptcache} delivers real savings but trusts a single threshold and never verifies a hit, so it is unsafe in exactly the look-alike cases that matter. The embedding and vector-search literature makes matching fast and semantic but says nothing about \emph{when a match should be believed}. The NLI literature provides the tool to verify equivalence but has not been integrated into a cache as a cheap-first, two-tier gate. And the routing literature reduces cost on the queries that miss the cache but is orthogonal to caching itself.

VeriGate closes this gap by unifying these threads into one gateway: an \emph{adaptive} threshold proposes candidates, a \emph{cheap-first two-tier verifier} disposes of false hits, a \emph{volatility} detector protects against staleness, a \emph{sub-linear vector index} keeps the hit path fast, and a \emph{cost-aware cascade} reduces spend on the residual traffic --- all wrapped in production-grade tenant isolation, budgeting, and observability. Table~\ref{tab:lit_comparison} contrasts VeriGate with the closest systems.

\begin{table}[H]
    \centering
    \resizebox{\textwidth}{!}{
    \begin{tabular}{|l|l|c|c|c|p{4.6cm}|}
    \hline
    \textbf{System} & \textbf{Core Idea} & \textbf{Adaptive Thr.} & \textbf{Hit Verification} & \textbf{Cost Routing} & \textbf{Primary Limitation / Gap} \tabularnewline \hline
    Exact cache (Redis) & String-key lookup & N/A & N/A & No & \raggedright Only helps on identical repeats; misses all paraphrases. \tabularnewline \hline
    GPTCache \cite{bang2023gptcache} & Fixed-threshold semantic cache & No & No & No & \raggedright False hits on look-alikes (number/entity/negation); brittle single threshold. \tabularnewline \hline
    FrugalGPT \cite{chen2023frugalgpt} & Cheap$\rightarrow$strong model cascade & N/A & N/A & Yes (learned) & \raggedright Reduces cost but has no caching; orthogonal to reuse. \tabularnewline \hline
    RouteLLM \cite{ong2024routellm} & Learned weak/strong router & N/A & N/A & Yes (learned) & \raggedright Needs preference-data training; no cache; no verification. \tabularnewline \hline
    \textbf{VeriGate (this work)} & \textbf{Adaptive self-verifying cache + cascade} & \textbf{Yes} & \textbf{Yes (2-tier)} & \textbf{Yes} & \raggedright \textbf{Curated benchmark; learned threshold/router left as future work.} \tabularnewline \hline
    \end{tabular}
    }
    \caption{Architectural comparison of VeriGate against the closest caching and routing systems.}
    \label{tab:lit_comparison}
\end{table}
